\section{Experiments}
\label{sec:exp}

\paragraph{Models.} \emph{Specialists}: STID~\citep{shao2022stid}, Graph WaveNet~\citep{wu2019graphwavenet} and
STAEformer~\citep{liu2023staeformer}, imported unchanged from their official repositories (pinned commits in the
supplement) and trained on each task's training split with the published hyper-parameters, loss, optimiser,
schedule and early stopping, on exactly the windows PWM-Dyn uses. The one code change is a compatibility fix in Graph
WaveNet, whose 2-D-kernel \texttt{Conv1d} layers are rejected by current PyTorch and are built as \texttt{Conv2d}, the
operation they computed. Each run had a wall-clock budget on one T4 (Table~\ref{tab:budget}); runs that hit it keep
their best-validation checkpoint and are marked. \emph{Foundation models}: TimesFM~2.5 (200M) and TimesFM~3.0 in its
multivariate mode, zero-shot, with a 512-step (365-day) context, run in a separate software environment.
\emph{References}: last value, seasonal naive and per-location climatology quantiles from the training split.
\emph{PWM-Dyn}: the configuration frozen by the ablation below, three fresh seeds (100--102), trained for the
optimisation-step budget it has on PEMS08 ($\approx$33.5k steps; 9--52 GPU-minutes per seed), conformally
calibrated on validation and scored once.

\begin{table*}[t]
\caption{Point accuracy: test MAE (RMSE below) on the official windows. Upper block: Track A, 12-step context, every model trained on the task's training split; lower block: Track B, foundation models zero-shot with 512 steps (365 days on rainfall). Bold: best MAE in the block. $^\dagger$\,stopped by its wall-clock budget before its authors' schedule ended (Table~\ref{tab:budget}). $^\ast$\,network never used during development. Rainfall PWM-Dyn entries are v0.4; rainfall has no specialist entries. Units: vehicles/5\,min (PeMS flow), mph (METR-LA, PEMS-BAY), mm/day (IMD).}
\label{tab:main}
\small
\setlength{\tabcolsep}{4pt}
\begin{tabular}{@{}lrrrrrrrr@{}}
\toprule
Model & PEMS08 & PEMS04 & PEMS03$^\ast$ & PEMS07$^\ast$ & METR-LA$^\ast$ & PEMS-BAY$^\ast$ & IMD-CI & IMD-WC \\
\midrule
Last value & 25.30 & 31.56 & 23.89 & 35.63 & 5.04 & 2.17 & 4.37 & 4.88 \\
 & {\scriptsize 38.66} & {\scriptsize 48.25} & {\scriptsize 37.31} & {\scriptsize 54.37} & {\scriptsize 10.98} & {\scriptsize 5.11} & {\scriptsize 13.55} & {\scriptsize 14.94} \\[1pt]
Climatology & 22.18 & 25.02 & 25.93 & 29.98 & 3.90 & 2.53 & 2.89 & 3.65 \\
 & {\scriptsize 40.77} & {\scriptsize 41.36} & {\scriptsize 47.80} & {\scriptsize 58.67} & {\scriptsize 8.41} & {\scriptsize 5.48} & {\scriptsize 10.70} & {\scriptsize 13.00} \\[1pt]
\midrule
STID & 14.18 & 18.49 & 15.50 & 19.62$^\dagger$ & 3.12 & 1.56$^\dagger$ & -- & -- \\
 & {\scriptsize 23.21} & {\scriptsize 30.05} & {\scriptsize 26.83} & {\scriptsize 32.72} & {\scriptsize 6.50} & {\scriptsize 3.60} &  &  \\[1pt]
Graph WaveNet & 14.77 & 19.37$^\dagger$ & 14.70$^\dagger$ & 23.39$^\dagger$ & 3.05$^\dagger$ & 1.57$^\dagger$ & -- & -- \\
 & {\scriptsize 23.59} & {\scriptsize 30.75} & {\scriptsize 25.17} & {\scriptsize 36.92} & {\scriptsize 6.15} & {\scriptsize 3.57} &  &  \\[1pt]
STAEformer & \textbf{13.93}$^\dagger$ & 19.21$^\dagger$ & 15.45$^\dagger$ & -- & 3.09$^\dagger$ & 1.59$^\dagger$ & -- & -- \\
 & {\scriptsize 23.17} & {\scriptsize 30.66} & {\scriptsize 26.03} &  & {\scriptsize 6.33} & {\scriptsize 3.63} &  &  \\[1pt]
PWM-Dyn (ours) & 14.12 & \textbf{18.46} & \textbf{14.34} & \textbf{19.56} & \textbf{2.92} & \textbf{1.52} & \textbf{2.84} & \textbf{3.41} \\
 & {\scriptsize 23.46} & {\scriptsize 30.77} & {\scriptsize 23.12} & {\scriptsize 32.60} & {\scriptsize 5.88} & {\scriptsize 3.45} & {\scriptsize 10.42} & {\scriptsize 12.01} \\[1pt]
\midrule
TimesFM 2.5 & 15.04 & 19.70 & -- & -- & -- & -- & 2.96 & 3.67 \\
 & {\scriptsize 24.09} & {\scriptsize 31.70} &  &  &  &  & {\scriptsize 10.58} & {\scriptsize 12.65} \\[1pt]
TimesFM 3.0 mv & 14.55 & 19.13 & \textbf{14.15} & 20.27 & 3.40 & 1.66 & 2.92 & 3.62 \\
 & {\scriptsize 23.57} & {\scriptsize 31.02} & {\scriptsize 22.70} & {\scriptsize 33.13} & {\scriptsize 6.97} & {\scriptsize 3.95} & {\scriptsize 10.65} & {\scriptsize 12.78} \\[1pt]
PWM-Dyn (ours), Track B & \textbf{14.18} & \textbf{18.46} & 14.34 & \textbf{19.56} & \textbf{2.92} & \textbf{1.52} & \textbf{2.84} & \textbf{3.41} \\
 & {\scriptsize 23.38} & {\scriptsize 30.77} & {\scriptsize 23.12} & {\scriptsize 32.60} & {\scriptsize 5.88} & {\scriptsize 3.45} & {\scriptsize 10.42} & {\scriptsize 12.01} \\[1pt]
\bottomrule
\end{tabular}
\end{table*}

\subsection{Point accuracy}
Table~\ref{tab:main} gives test MAE and RMSE. Under the standard 12-step context (upper block) the PWM-Dyn ensemble
has the lowest MAE on PEMS04, PEMS03, PEMS07, METR-LA and PEMS-BAY and is second on PEMS08 (14.12) behind STAEformer
(13.93), which reached only 25 of its scheduled epochs; STAEformer's published PEMS08 result, 13.46 on all test
windows, is better still. On the four networks never used during development the margin over the best trained
opponent is 2.5\% (PEMS03, Graph WaveNet), 0.3\% (PEMS07, STID), 4.2\% (METR-LA, Graph WaveNet) and 2.7\% (PEMS-BAY,
STID). STID, which completed its authors' schedule on four networks, is the closest rival on PEMS08, PEMS04, PEMS07 and
PEMS-BAY and essentially ties PWM-Dyn on PEMS08 and PEMS04. Part of PWM-Dyn's margin is ensembling: its three single
seeds average 14.33, 18.62 and 19.88 on PEMS08, PEMS04 and PEMS07, slightly worse than single-seed STID (14.18, 18.49,
19.62), while the ensemble costs about ten times STID's training time. In the open-context block, TimesFM~3.0 run zero-shot beats
every trained model on PEMS03 (14.15) and is ahead of or within 4\% of PWM-Dyn on four of six traffic networks;
PWM-Dyn's long-context variant was preferred on validation only for PEMS08. RMSE orders the models similarly, with
exceptions: all three specialists on PEMS04, and STID and STAEformer on PEMS08, have lower RMSE than PWM-Dyn, i.e.\
smaller errors on the largest deviations.


\subsection{Is the difference larger than noise?}
\begin{table}[t]
\caption{Paired tests of PWM-Dyn against each model on the same windows: relative MAE difference (negative = PWM-Dyn better) and Diebold--Mariano $p$-value; in brackets the 95\% day-block bootstrap interval. Track A for specialists, Track B for TimesFM. $^\dagger$\,opponent budget-capped.}
\label{tab:sig}
\scriptsize
\setlength{\tabcolsep}{2.5pt}
\begin{tabular}{@{}lcccc@{}}
\toprule
Task & STID & GWNet & STAEf. & TimesFM 3.0 \\
\midrule
PEMS08 & \shortstack{$-0.4\%$, $p=0.12$\\{}[-1.0, +0.1]} & \shortstack{$-4.4\%$, $p<10^{-3}$\\{}[-4.9, -4.0]} & \shortstack{$+1.4\%$, $p<10^{-3}$$^\dagger$\\{}[+0.7, +1.9]} & \shortstack{$-2.5\%$, $p=0.06$\\{}[-5.4, -0.3]} \\[2pt]
PEMS04 & \shortstack{$-0.2\%$, $p=0.82$\\{}[-2.0, +1.1]} & \shortstack{$-4.7\%$, $p<10^{-3}$$^\dagger$\\{}[-5.3, -4.3]} & \shortstack{$-3.9\%$, $p<10^{-3}$$^\dagger$\\{}[-4.7, -3.5]} & \shortstack{$-3.5\%$, $p=0.03$\\{}[-7.0, -0.8]} \\[2pt]
PEMS03$^\ast$ & \shortstack{$-7.5\%$, $p=0.02$\\{}[-14.1, -2.3]} & \shortstack{$-2.5\%$, $p<10^{-3}$$^\dagger$\\{}[-3.3, -1.2]} & \shortstack{$-7.2\%$, $p<10^{-3}$$^\dagger$\\{}[-10.8, -4.6]} & \shortstack{$+1.3\%$, $p=0.48$\\{}[-1.8, +5.3]} \\[2pt]
PEMS07$^\ast$ & \shortstack{$-0.3\%$, $p=0.03$$^\dagger$\\{}[-0.5, +0.0]} & \shortstack{$-16.4\%$, $p<10^{-3}$$^\dagger$\\{}[-17.5, -15.4]} & -- & \shortstack{$-3.5\%$, $p<10^{-3}$\\{}[-5.4, -1.6]} \\[2pt]
METR-LA$^\ast$ & \shortstack{$-6.6\%$, $p<10^{-3}$\\{}[-7.6, -5.5]} & \shortstack{$-4.2\%$, $p<10^{-3}$$^\dagger$\\{}[-5.0, -3.3]} & \shortstack{$-5.5\%$, $p<10^{-3}$$^\dagger$\\{}[-6.2, -4.8]} & \shortstack{$-14.2\%$, $p<10^{-3}$\\{}[-16.3, -11.5]} \\[2pt]
PEMS-BAY$^\ast$ & \shortstack{$-2.7\%$, $p<10^{-3}$$^\dagger$\\{}[-3.6, -1.9]} & \shortstack{$-3.2\%$, $p<10^{-3}$$^\dagger$\\{}[-3.9, -2.2]} & \shortstack{$-4.7\%$, $p<10^{-3}$$^\dagger$\\{}[-5.2, -4.2]} & \shortstack{$-8.5\%$, $p<10^{-3}$\\{}[-11.4, -5.2]} \\[2pt]
\bottomrule
\end{tabular}
\end{table}
Table~\ref{tab:sig} tests each comparison on the per-window error differences. We call a difference significant only
when the Diebold--Mariano test rejects at $p<0.05$ \emph{and} the day-block bootstrap interval excludes zero. Against
STID, PWM-Dyn's advantage is significant on METR-LA, PEMS-BAY and PEMS03 ($p\le0.023$; on PEMS03 the interval is
wide because the test period contains Thanksgiving week). On PEMS07 it is not ($-0.3\%$, $p=0.03$, interval
$[-0.5,+0.0]$), and on PEMS08 and PEMS04 the two models are statistically tied ($p=0.12$ and $0.82$). STAEformer's lead
on PEMS08 is significant. Against TimesFM~3.0 the advantage is significant on PEMS04, PEMS07, METR-LA and PEMS-BAY; on
PEMS08 the interval excludes zero but $p=0.06$; on PEMS03 TimesFM is ahead by 1.3\%, not significantly. The standard
deviation of single-seed test MAE is 0.002--0.22 (Table~\ref{tab:abl}, right), small next to the margins over Graph
WaveNet, STAEformer and TimesFM but as large as the gaps to STID on PEMS08, PEMS04 and PEMS07; the ensemble beats each
of its seeds on every network.


\subsection{Uncertainty and dead sensors}
\begin{table}[t]
\caption{Probabilistic quality (Track B, IS80 with 80\% coverage) and dead sensors (Track R, MAE on the 30\% of locations with no test-period readings). Specialists without a missing-data mechanism are fed zeros, their repositories' missing-value convention.}
\label{tab:prob}
\small
\setlength{\tabcolsep}{3pt}
\begin{tabular}{@{}lrrrrr@{}}
\toprule
 & \multicolumn{3}{c}{IS80 (coverage \%)} & \multicolumn{2}{c}{Track R MAE} \\
\cmidrule(lr){2-4}\cmidrule(l){5-6}
Task & PWM-Dyn & TimesFM 3.0 & Climat. & PWM-Dyn & best other \\
\midrule
PEMS08 & 66.94 (80) & 67.84 (81) & 117.83 (71) & \textbf{19.25} & 21.78\,{\scriptsize(clim.)} \\
PEMS04 & 86.23 (80) & 88.42 (81) & 127.81 (72) & \textbf{25.33} & 27.39\,{\scriptsize(clim.)} \\
PEMS03$^\ast$ & 67.14 (78) & 65.50 (81) & 154.10 (61) & \textbf{17.71} & 23.31\,{\scriptsize(clim.)} \\
PEMS07$^\ast$ & 93.99 (80) & 95.32 (81) & 159.70 (76) & 30.95 & \textbf{29.10}\,{\scriptsize(clim.)} \\
METR-LA$^\ast$ & 14.34 (81) & 17.37 (81) & 19.42 (75) & \textbf{3.36} & 4.07\,{\scriptsize(clim.)} \\
PEMS-BAY$^\ast$ & 7.54 (79) & 8.29 (81) & 13.19 (75) & \textbf{2.31} & 2.71\,{\scriptsize(clim.)} \\
IMD-CI & 18.01 (85) & 19.79 (86) & 18.21 (91) & \textbf{2.84} & 2.88\,{\scriptsize(clim.)} \\
IMD-WC & 20.27 (83) & 21.85 (87) & 21.09 (91) & \textbf{3.43} & 3.60\,{\scriptsize(clim.)} \\
\bottomrule
\end{tabular}
\end{table}
\textbf{Intervals.} After conformal calibration, PWM-Dyn's 80\% intervals cover 78--85\% of test values and have a
lower interval score than TimesFM~3.0's on seven of eight tasks (Table~\ref{tab:prob}; the two rainfall entries are
from the previous model version); the exception is again PEMS03. TimesFM's intervals slightly over-cover (81--87\%)
everywhere; climatology intervals are wide and under-cover on traffic. Specialists give point forecasts only.
\textbf{Dead sensors.} With 30\% of locations dark throughout the test period, PWM-Dyn's MAE on those locations rises
by 15--58\% over its normal accuracy on traffic (and hardly at all on rainfall) but stays below per-sensor climatology
on seven of eight tasks; on PEMS07 it is 6\% worse than climatology. On PEMS08 and PEMS04 the previous version was
also slightly better on dark sensors (18.70 and 24.99 against 19.25 and 25.33); our ablations do not isolate why. The
standard-context model's lack of a climatology fallback is one candidate on PEMS04 and the new networks, but it does
not explain PEMS08, where both versions keep it. The three specialists, fed zeros for the missing readings as their
repositories do, produce errors of 5--310 (appendix): forecasting dark sensors needs an explicit mechanism, not just a
graph.


\subsection{What cost v0.4 its accuracy}
\begin{table}[t]
\caption{Left: what cost v0.4 accuracy. Validation MAE (selection criterion, seeds 0--1) and test MAE of each configuration's ensemble; the frozen choice is marked. Right: test MAE of the three fresh seeds of the frozen model and of their ensemble (Track A). The v0.4.1 row is its own three-seed run (seeds 0--2).}
\label{tab:abl}
\scriptsize
\setlength{\tabcolsep}{2.5pt}
\begin{tabular}{@{}lrrrr@{}}
\toprule
Configuration & \multicolumn{2}{c}{Val.\ MAE} & \multicolumn{2}{c}{Test MAE} \\
\cmidrule(lr){2-3}\cmidrule(l){4-5}
 & P08 & P04 & P08 & P04 \\
\midrule
clim., 100 ep., select every 4th (v0.4.1) & 14.52 & 18.84 & 14.50 & 19.02 \\
clim., 100 ep., select every window & 14.45 & 18.84 & 14.45 & 19.06 \\
no clim., 60 ep., select every window & 14.60 & 18.61 & 14.23 & 18.47 \\
no clim., 100 ep., select every window (frozen) & 14.61 & 18.39 & 14.14 & 18.42 \\
\bottomrule
\end{tabular}
\vspace{4pt}
\begin{tabular}{@{}lrrr@{}}
\toprule
Task & Seeds (test MAE) & Mean $\pm$ SD & Ens. \\
\midrule
PEMS08 & 14.22 / 14.24 / 14.52 & 14.33 $\pm$ 0.17 & 14.12 \\
PEMS04 & 18.52 / 18.70 / 18.65 & 18.62 $\pm$ 0.09 & 18.46 \\
PEMS03$^\ast$ & 14.45 / 14.45 / 14.65 & 14.52 $\pm$ 0.12 & 14.34 \\
PEMS07$^\ast$ & 20.12 / 19.69 / 19.83 & 19.88 $\pm$ 0.22 & 19.56 \\
METR-LA$^\ast$ & 2.960 / 2.959 / 2.956 & 2.958 $\pm$ 0.002 & 2.919 \\
PEMS-BAY$^\ast$ & 1.544 / 1.536 / 1.535 & 1.538 $\pm$ 0.005 & 1.519 \\
\bottomrule
\end{tabular}
\end{table}
An earlier version of PWM-Dyn added a per-location climatology input, and its test MAE under the 12-step context was
2--4\% worse than the version before it. Table~\ref{tab:abl} separates the changes involved on validation. On PEMS04
the climatology input accounts for most of the loss: with it, validation MAE is 18.84 whichever checkpoint rule is
used; without it, 18.39--18.61, with training length (60 vs 100 epochs) moving it by 0.22. On PEMS08 the effect is
mixed: the climatology input improves validation MAE (14.45 vs 14.61) but not test MAE. The configuration with the
lowest validation MAE averaged over both networks (no climatology input, 100 epochs, checkpoint chosen on every
validation window) was frozen before any test scoring; its test MAE, reported for the record, is also the lowest of
the four. The climatology input is not useless: it gives dark sensors a sensible baseline and remains in the
long-context variant.


\subsection{Cost}
PWM-Dyn has 0.32--0.41M parameters, against 0.12--0.14M (STID), 0.31--0.32M (Graph WaveNet), 1.2--1.4M (STAEformer)
and 200M (TimesFM~2.5). One seed trains in 9 (PEMS08), 9 (METR-LA), 13 (PEMS04), 14 (PEMS-BAY), 17 (PEMS03) and 52
(PEMS07) minutes on a T4, so the three-seed ensembles cost 0.4--2.6 GPU-hours per network; TimesFM~3.0 needed
0.6--2.0 GPU-hours per network for inference alone at our window counts. The study reported here ran in two free Kaggle sessions on 2$\times$T4 (4.6 and 6.5 hours, at most
22 T4-hours including TimesFM inference), plus the earlier runs that produced the rainfall entries.

